\section{Witt--Bose--Shrikhande local actions}
We use the standard nonsplit-torus construction of the
Witt--Bose--Shrikhande space, attributed in source Section~3.4 to the earlier
work cited there. Its full automorphism group
$\operatorname{P\Gamma L}_2(q)$ is the retained classification input from
BDD, as used by the source; the local-action argument below is given in
full. In particular, the full-group identification is not inferred from
an exhibited subgroup. Since $q$ is even,
$\operatorname{PSL}_2(q)=\operatorname{PGL}_2(q)=\operatorname{SL}_2(q)$.

\begin{proposition}[Torus--involution incidence]\label{wbs-model}
Let $q=2^e$, $e\geq3$, and $G=\operatorname{SL}_2(q)$. Let the points be
normalizers $D=N_G(T)$ of nonsplit tori $T$ of order $q+1$, and let the lines
be the involutions $s$ of $G$. Declare $D$ incident with $s$ when $s\in D$.
This is a flag-transitive linear space with
\[
 v=q(q-1)/2,\qquad b=q^2-1,\qquad k=q/2,\qquad r=q+1.
\]
A point stabilizer in $G$ is $D$ of order $2(q+1)$, a line stabilizer is
$U=C_G(s)$ of order $q$, and a flag stabilizer is $\langle s\rangle$.
\end{proposition}
\begin{proof}
Identify the two-dimensional vector space with $\F_{q^2}$. Multiplication
by its norm-one elements gives the cyclic group $T$. The map $z\mapsto z^q$
inverts $T$ and, in characteristic two, has determinant one. The normalizer
is therefore the dihedral group $D$ of order $2(q+1)$: its elements either
fix or exchange the two common eigenlines over $\F_{q^2}$. Conversely these
are all possibilities for a normalizer. Every nonidentity element of $T$
has distinct eigenvalues outside $\F_q$, since
$\gcd(q-1,q+1)=1$, and its centralizer in $G$ is exactly $T$.
All nonsplit tori are conjugate: conjugate their two-dimensional field
embeddings, and adjust a conjugating matrix's determinant by a centralizing
field multiplier, whose norm ranges over $\F_q^\times$. Also $D$ is
self-normalizing, since its odd cyclic subgroup is characteristic.
This proves the point count.

Every involution is conjugate to
$s=\left(\begin{smallmatrix}1&1\\0&1\end{smallmatrix}\right)$.
Its centralizer consists exactly of
$u_a=\left(\begin{smallmatrix}1&a\\0&1\end{smallmatrix}\right)$,
$a\in\F_q$, as direct matrix multiplication shows. Thus $U$ is elementary
abelian of order $q$ and there are $q^2-1$ involutions. The $q+1$ reflections
of $D$ form one $T$-orbit, because squaring is bijective on $T$.
Consequently $G$ is flag-transitive and the flag stabilizer is
$C_D(s)=\langle s\rangle$. This proves $r=q+1$ and $k=q/2$.

Two distinct point normalizers cannot contain two distinct common
involutions: their nontrivial product lies in the rotation torus of each,
and its unique centralizer identifies those tori and hence the normalizers.
Thus two points lie on at most one line. Moreover
\[
 b\binom{k}{2}=\frac{(q^2-1)q(q-2)}8
             =\binom{q(q-1)/2}{2}.
\]
Every pair therefore lies on exactly one line. The parameter bounds give
at least two points on every line and two lines on every point.
\end{proof}

\begin{theorem}[Uniform local exclusion and a corrected source sentence]
\label{wbs-exclusion}
For $e\geq3$, no flag-transitive subgroup of the full group
$A=\operatorname{P\Gamma L}_2(2^e)$ of $W(2^e)$ is residually weakly
primitive. Nevertheless the full group of $W(8)$ is primitive on the points
of each line: its induced line action is $\operatorname{AGL}_1(4)\cong A_4$.
Thus the final WBS exclusion of source Section~3.4 is valid, but its stronger
assertion that the full group is never locally point-primitive is false.
\end{theorem}
\begin{proof}
First consider the pencil at $D$. Its lines are the $q+1$ involutions in
$D$. The rotation subgroup $T$ acts regularly on them. It is normal in
$A_D$, since it is characteristic in $D$ and $G\lhd A$. If $q+1$ is
composite, a proper nontrivial characteristic subgroup of the cyclic $T$
gives a nontrivial $A_D$-invariant partition into its orbits. Thus pencil
primitivity forces $2^e+1$ to be prime. Writing $e=2^a m$ with $m$ odd
shows that $m=1$: for $m>1$, $2^{2^a}+1$ is a proper divisor. Hence
$e$ is a power of two, and $e\geq4$ in the present range.

Now consider the line represented by the displayed involution $s=u_1$.
Its stabilizer is the centralizer $A_s$. Each element of $A$ is a matrix
part followed by a field automorphism. Field automorphisms fix $s$;
the matrix centralizer calculation above gives
\[
 A_s=U\rtimes\langle\sigma\rangle,
 \qquad \sigma(a)=a^2,\qquad |\langle\sigma\rangle|=e.
\]
Fix a point $D$ on this line and put $B=A_s\cap A_D$.
Since $G$ is flag-transitive, every coset of $G$ in $A$ has a representative
fixing this flag. Hence $B/(B\cap G)\cong A/G\cong C_e$, whereas
$B\cap U=B\cap G=\langle s\rangle$. Also $A_s=UB$.
The line points can consequently be identified with the additive quotient
\[
 V=U/\langle s\rangle\cong\F_{2^e}/\F_2,
\]
with $V$ acting regularly. Conjugation by the point stabilizer acts on this
quotient through $\langle\sigma\rangle$; translation parts have no effect
on that action. The affine block criterion in
Proposition~\ref{translation} says that line primitivity is equivalent to
irreducibility of this stabilizer on $V$.

For the only remaining pencil possibility, $e$ is a power of two. The
stabilizer image on $V$ is therefore a $2$-group. A $2$-group acting on a
nonzero vector space over $\F_2$ has a nonzero fixed vector: all its nonfixed
orbits have even size, so the number of fixed vectors is even and includes
zero. Since $\dim_{\F_2}V=e-1>1$, a fixed one-dimensional subspace is a
proper invariant subspace. The line action is imprimitive. Thus the full
group fails at least one of the two required local actions for every $e$.
A primitive local action of a subgroup would make the corresponding full
group action primitive by Lemma~\ref{coarsening}; no subgroup can evade
this obstruction. This proves the uniform exclusion.

For the stronger source assertion, set $e=3$. On $V=\F_8/\F_2$, a class
fixed by $\sigma$ has $x^2+x\in\F_2$. The value zero gives the zero class;
the value one would put a root of $X^2+X+1$ in $\F_8$, impossible because
its roots have multiplicative order three whereas $|\F_8^\times|=7$.
The order-three action on the two-dimensional $V$ is therefore irreducible.
It is transitive on the three nonzero vectors. The induced line group is
$V\rtimes C_3$ of order twelve, proving the claimed primitive action.
Its pencil remains imprimitive since $q+1=9$. The short accompanying
$q=8$ calculation validates these exact stabilizers but is not a premise
of the written proof.
\end{proof}
